\scanpage{131}
\uline{Proposition} Soit $Q\in\mathbb Q^*$, et soit $\lambda\in\mathbb C$ un nombre algébrique.
Conditions équivalentes

(a) Pour tout conjugué $\lambda_\alpha$ de $\lambda$, on a $\lambda_\alpha\bar\lambda_\alpha=Q$ i.e. $|\lambda_\alpha|^2=\unsure{Q}$

(b) \deleted{\unsure{Il est équivalent}} Il existe un corps \inserted{de nbres} \uline{totalement réel} $K$, tel que
\deleted{\unsure{son s}} $K(\lambda)$ soit égal à $K$ ou une extension \uline{quadratique} purement
imaginaire de $K$, et on a $\lambda\bar\lambda=Q$ i.e. $|\lambda|^2=Q$.

(b') Comme (b), avec $K=\mathbb Q(a)\deleted{\unsure{=K\cap\mathbb R}}$, où $a=\lambda+\bar\lambda=\lambda+\frac Q\lambda$.

(b'') Il existe un nombre algébrique \deleted{\unsure{et|$a$}} \uline{tot} \uline{réel} $a$,
avec \deleted{$\unsure{a^2-4Q}$ nul} $4Q-a^2$ nul ou totalement
positif, tel que $\lambda$ satisfasse une équation
\[
\lambda^2-a\lambda+Q=0
\]

(c) Soit $\mathbb K$ l'extension alg. tot. réelle maximale de $\mathbb Q$, et
$\mathbb L$ l'unique extension quadratique de $\mathbb K$ galoisienne sur $\mathbb Q$.
On a $\lambda\in\mathbb L$, et \deleted{$\unsure{|\lambda|}$} $\lambda\bar\lambda=Q$ i.e. $|\lambda|^2=Q$.

\uline{Dém}

a) $\Longrightarrow$ b'') Posons $a=\lambda+\bar\lambda=\lambda+\frac Q\lambda$. Je dis que $a$ \uline{est tot. réel}.
En effet, les conjugués de $a$ sont les $\deleted{\unsure{\sigma}}a_\alpha=\deleted{\unsure{\sigma}}\lambda_\alpha+\frac{Q}{\deleted{\unsure{\sigma}}\lambda_\alpha}$, or
par hypothèse $\lambda_\alpha\bar\lambda_\alpha=Q$ donc $a_\alpha=\lambda_\alpha+\bar\lambda_\alpha\in\mathbb R$. D'autre part,
$\lambda$ est solution de l'équation
\[
(T-\lambda)(T-\bar\lambda)=0\qquad\text{i.e.}\qquad T^2-aT+Q=0,
\]
et par suite $\lambda_\alpha$ est solution de l'équation $T^2-a_\alpha T+Q_{\deleted{\unsure{q}}}=0$,
\deleted{de plus} l'autre solution étant $\frac Q{\lambda_\alpha}$\inserted{(qui est donc égale à $\bar\lambda_\alpha$)}. \deleted{\unsure{Donc}} Comme
$\lambda_\alpha\bar\lambda_\alpha=Q$, les deux solutions sont $\lambda_\alpha$ et $\bar\lambda_\alpha$, donc complexes
conjuguées. Cela montre que $a_\alpha^2-4Q\le0$, \deleted{\unsure{donc|d'où}} \inserted{i.e.}
$4Q-a_\alpha^2\ge0$, donc $4Q-a^2$ est tot positif.

b'') $\Longrightarrow$ b') Il est évident que $\lambda$ est de degré au plus
deux sur \inserted{le corps tot. réel} $\mathbb Q(a)$, \deleted{De plus \unsure{comme} $\unsure{a=\lambda+\frac Q\lambda}$} \inserted{et $\mathbb Q(\lambda)\supset\mathbb Q(a)$ [car $a=\lambda+\frac Q\lambda$].}
\deleted{\unsure{ce qui}} Comme $a^2-4Q\le0$, les solutions de l'équation
$T^2-aT+Q=0$ sont $\lambda$ et $\bar\lambda$\inserted{(cf lemme plus bas)}, et on a $\lambda\bar\lambda=Q$, $\lambda\unsure{+|-}\bar\lambda=a$.
% archive sheet number: 130
